% =============================================================================
% APPENDIX MMM: METHOD-POLICY-COHERENCE
% Cross-Domain Spillovers and Multi-Domain Policy Coherence Framework
% =============================================================================
%
% Code: MMM
% Category: METHOD-POLICY-COHERENCE (Cross-domain integration)
% Language: English
% Version: 1.0 (January 2026)
%
% Purpose: Provide systematic framework for analyzing how intervention
% portfolios in one policy domain (Health, Pension, Environment, Labor)
% affect outcomes in other domains. Includes spillover matrices, decision
% trees for conflict detection, coherence scoring, and case studies of
% policy conflicts and complementarities.
%
% For: Government policy coordinators, multi-sector implementation teams,
% politicians designing comprehensive policy agendas who need to understand
% and manage policy interactions.
%
% Dependencies: Chapter 22 (Policy Applications), Appendix LLL (domain templates)
%
% Cross-references: Appendix HHH (coherence criteria), KKK (dynamics)
%
% =============================================================================

\newpage
\section*{Appendix MMM: METHOD-POLICY-COHERENCE}
\label{app:method-policy-coherence}
\addcontentsline{toc}{section}{Appendix MMM: METHOD-POLICY-COHERENCE}

% =============================================================================
% HEADER
% =============================================================================

\begin{tcolorbox}[colback=blue!5!white, colframe=blue!75!black, title=Quick Reference]
\textbf{Appendix:} MMM (METHOD)

\small
\textbf{Purpose:} Analyze how policy portfolios in one domain affect other domains. Detect conflicts, identify synergies, optimize cross-domain coherence.

\textbf{3-Step Process:}
\begin{enumerate}[nosep]
\item Map your current policies across all domains (Section 1)
\item Compute spillover effects using matrices (Sections 2-3)
\item Optimize coherence score (>0.25 acceptable, >0.35 good) and redesign conflicting policies (Section 4)
\end{enumerate}

\textbf{Measurement:} Review spillovers quarterly alongside domain-level KPIs
\end{tcolorbox}



\begin{abstract}
This appendix provides systematic analysis and documentation relevant to the
Evidence-Based Framework (EBF). It integrates with the 10C CORE architecture
and provides validated findings for practical application.
\end{abstract}

\begin{tcolorbox}[
    colback=orange!5!white,
    colframe=orange!75!black,
    title={\textbf{Appendix Scope: Ziel / In-Scope / Out-of-Scope / Constraints}},
    fonttitle=\bfseries
]

\textbf{Ziel (Objective):}
\begin{itemize}[nosep]
    \item Provide systematic analysis for METHOD integration
\end{itemize}

\textbf{In-Scope:}
\begin{itemize}[nosep]
    \item Core analysis and findings
    \item Framework integration points
\end{itemize}

\textbf{Out-of-Scope:}
\begin{itemize}[nosep]
    \item Detailed empirical replication $\rightarrow$ METHOD appendices
\end{itemize}

\textbf{Constraints:}
\begin{itemize}[nosep]
    \item Analysis bounded by available data
\end{itemize}

\end{tcolorbox}

---

% =============================================================================
% SECTION 1: MAPPING POLICIES ACROSS DOMAINS
% =============================================================================

\subsection*{Section 1: Multi-Domain Policy Inventory}

Before analyzing spillovers, document all active policies:

\begin{center}
\small
\begin{tabular}{@{}p{2cm}p{2.5cm}p{2.5cm}p{2.5cm}@{}}
\toprule
\textbf{Domain} & \textbf{Key Policies} & \textbf{Target Segment} & \textbf{Phase Focus} \\
\midrule
\textbf{HEALTH} & [1] \_\_\_\_\_\_\_\_\_\_\_\_ & [Segment] & [Phase] \\
 & [2] \_\_\_\_\_\_\_\_\_\_\_\_ & [Segment] & [Phase] \\
\midrule
\textbf{PENSION} & [1] \_\_\_\_\_\_\_\_\_\_\_\_ & [Segment] & [Phase] \\
 & [2] \_\_\_\_\_\_\_\_\_\_\_\_ & [Segment] & [Phase] \\
\midrule
\textbf{ENVIRONMENT} & [1] \_\_\_\_\_\_\_\_\_\_\_\_ & [Segment] & [Phase] \\
 & [2] \_\_\_\_\_\_\_\_\_\_\_\_ & [Segment] & [Phase] \\
\midrule
\textbf{LABOR} & [1] \_\_\_\_\_\_\_\_\_\_\_\_ & [Segment] & [Phase] \\
 & [2] \_\_\_\_\_\_\_\_\_\_\_\_ & [Segment] & [Phase] \\
\midrule
\end{tabular}
\end{center}

---

% =============================================================================
% SECTION 2: CROSS-DOMAIN SPILLOVER MATRICES
% =============================================================================

\subsection*{Section 2: Complete Spillover Matrix (4×4 Domain Pairs)}

\textbf{Instructions:} For each cell, estimate γ between domains (range: -1 to +1)
\begin{itemize}[nosep]
\item γ > +0.3 = Strong positive spillover (policies reinforce)
\item γ = 0 to +0.3 = Weak positive spillover (neutral, slight reinforcement)
\item γ = 0 to -0.3 = Weak negative spillover (neutral, slight interference)
\item γ < -0.3 = Strong negative spillover (policies interfere) → \textbf{CONFLICT RISK}
\end{itemize}

\subsubsection*{2.1 Master Spillover Matrix}

\begin{center}
\renewcommand{\arraystretch}{1.8}
\begin{tabular}{@{}c|cccc@{}}
\toprule
\textbf{From→To} & \textbf{Health} & \textbf{Pension} & \textbf{Environment} & \textbf{Labor} \\
\midrule
\textbf{Health} & — & \_\_\_\_ & \_\_\_\_ & \_\_\_\_ \\
\textbf{Pension} & \_\_\_\_ & — & \_\_\_\_ & \_\_\_\_ \\
\textbf{Environment} & \_\_\_\_ & \_\_\_\_ & — & \_\_\_\_ \\
\textbf{Labor} & \_\_\_\_ & \_\_\_\_ & \_\_\_\_ & — \\
\midrule
\end{tabular}
\end{center}

\textbf{Interpretation:}
\begin{itemize}[nosep]
\item Read row→column to find spillover effect (e.g., Health→Environment)
\item γ symmetric approximation: If Health→Env = +0.4, assume Env→Health ≈ +0.4
\item Fill in estimates based on your specific policies
\end{itemize}

\subsubsection*{2.2 Typical Spillover Patterns (Reference Table)}

\begin{center}
\small
\begin{tabular}{@{}llcl@{}}
\toprule
\textbf{Domain Pair} & \textbf{Typical Spillover} & \textbf{γ} & \textbf{Mechanism} \\
\midrule
Health ↔ Environment & Aligned (personal + planet care) & +0.3–+0.5 & Both appeal to identity \\
Health ↔ Pension & Conflicting (invest now vs. save now) & 0.0–+0.1 & Different urgencies \\
Health ↔ Labor & Neutral to positive (healthy workers) & +0.1–+0.3 & Complementary \\
Pension ↔ Labor & Strong positive (work-life link) & +0.5–+0.7 & Direct connection \\
Pension ↔ Environment & Weak (different timeframes) & +0.0–+0.2 & Competing narratives \\
Environment ↔ Labor & Green jobs narrative & +0.1–+0.3 & Potential synergy \\
\midrule
\end{tabular}
\end{center}

---

% =============================================================================
% SECTION 3: SPILLOVER ANALYSIS BY DOMAIN PAIR
% =============================================================================

\subsection*{Section 3: Detailed Spillover Analysis}

\subsubsection*{3.1 Health ↔ Pension Spillover}

\textbf{Your Policies:}
\begin{center}
\begin{tabular}{@{}ll@{}}
\toprule
\textbf{Domain} & \textbf{Policy} \\
\midrule
Health & [INSERT: e.g., ``Promote exercise and healthy diet''] \\
Pension & [INSERT: e.g., ``Flexible retirement age 60-67''] \\
\midrule
\end{tabular}
\end{center}

\textbf{Spillover Analysis:}

\begin{quote}
\textbf{Congruence:} Do policies send the same message? \\
Health says: ``Invest in your health \textit{now}'' \\
Pension says: ``Work \textit{longer} to fund retirement'' \\
Alignment Score: \_\_\_/10 (1=conflicting, 10=perfectly aligned)

\textbf{Capacity:} Do they compete for implementation resources? \\
Health implementation: [Your department] \\
Pension implementation: [Your department] \\
Resource competition: [Low / Medium / High] \\

\textbf{Behavioral Coherence:} Same or different population segments? \\
Health targets: [Segment] \\
Pension targets: [Segment] \\
Overlap: \_\_\_\%
\end{quote}

\textbf{Spillover Estimate:} γ = \_\_\_\_ (expected range: 0.0 to +0.1)

\subsubsection*{3.2 Health ↔ Environment Spillover}

\textbf{Your Policies:}
\begin{center}
\begin{tabular}{@{}ll@{}}
\toprule
\textbf{Domain} & \textbf{Policy} \\
\midrule
Health & [INSERT: e.g., ``Promote physical activity, reduce sedentary behavior''] \\
Environment & [INSERT: e.g., ``Reduce energy consumption 20\%''] \\
\midrule
\end{tabular}
\end{center}

\textbf{Spillover Analysis:}

\begin{quote}
\textbf{Congruence:} Both appeal to personal agency and care \\
Health: ``Take care of your health'' \\
Environment: ``Take care of the planet'' \\
Message Synergy: [INSERT assessment]

\textbf{Capacity:} Largely independent implementation \\
Resource competition: [Low / Medium / High] \\

\textbf{Behavioral Coherence:} Both target autonomy-seeking + social-oriented \\
Overlap: [High / Medium / Low]
\end{quote}

\textbf{Spillover Estimate:} γ = \_\_\_\_ (expected range: +0.3 to +0.5)

\subsubsection*{3.3 Pension ↔ Labor Market Spillover}

\textbf{Your Policies:}
\begin{center}
\begin{tabular}{@{}ll@{}}
\toprule
\textbf{Domain} & \textbf{Policy} \\
\midrule
Pension & [INSERT: e.g., ``Flexible retirement 60-67 years''] \\
Labor & [INSERT: e.g., ``Skills training for workers 55+''] \\
\midrule
\end{tabular}
\end{center}

\textbf{Spillover Analysis:}

\begin{quote}
\textbf{Congruence:} Both directly affect work-life decisions \\
Pension message: ``Flexibility to retire when ready'' \\
Labor message: ``Reskill to extend your career'' \\
Message Coherence: [Strongly aligned / Somewhat conflicting / Neutral]

\textbf{Capacity:} Both require employer coordination \\
Resource competition: [High / Medium / Low] \\

\textbf{Behavioral Coherence:} Target: everyone 55+ \\
Overlap: Very high (95\%+)
\end{quote}

\textbf{Spillover Estimate:} γ = \_\_\_\_ (expected range: +0.5 to +0.7)

\subsubsection*{3.4 Environment ↔ Labor Spillover}

\textbf{Your Policies:}
\begin{center}
\begin{tabular}{@{}ll@{}}
\toprule
\textbf{Domain} & \textbf{Policy} \\
\midrule
Environment & [INSERT: e.g., ``Transition to renewable energy''] \\
Labor & [INSERT: e.g., ``Reskill fossil fuel workers''] \\
\midrule
\end{tabular}
\end{center}

\textbf{Spillover Analysis:}

\begin{quote}
\textbf{Congruence:} ``Green jobs'' narrative \\
Environment message: ``Sustainable future'' \\
Labor message: ``New opportunities in green sector'' \\
Message Alignment: [Strong / Moderate / Weak]

\textbf{Capacity:} Different departments but linked \\
Resource competition: [Low / Medium / High] \\

\textbf{Behavioral Coherence:} Target: displaced workers + green job seekers \\
Overlap: [High / Medium / Low]
\end{quote}

\textbf{Spillover Estimate:} γ = \_\_\_\_ (expected range: +0.1 to +0.3)

---

% =============================================================================
% SECTION 4: COHERENCE SCORING AND OPTIMIZATION
% =============================================================================

\subsection*{Section 4: Policy Coherence Score and Recommendations}

\subsubsection*{4.1 Computing Your Coherence Score}

\textbf{Formula:}
$$\text{Coherence Score} = \frac{\sum_{i < j} \gamma_{ij}}{\text{\# of domain pairs}}$$

For 4 domains, there are $\binom{4}{2} = 6$ pairs.

\textbf{Your Calculation:}

\begin{center}
\begin{tabular}{@{}llc@{}}
\toprule
\textbf{Domain Pair} & \textbf{Your γ Estimate} & \textbf{Typical Range} \\
\midrule
Health ↔ Pension & \_\_\_\_ & 0.0–0.1 \\
Health ↔ Environment & \_\_\_\_ & +0.3–+0.5 \\
Health ↔ Labor & \_\_\_\_ & +0.1–+0.3 \\
Pension ↔ Environment & \_\_\_\_ & +0.0–+0.2 \\
Pension ↔ Labor & \_\_\_\_ & +0.5–+0.7 \\
Environment ↔ Labor & \_\_\_\_ & +0.1–+0.3 \\
\midrule
\textbf{Sum} & \textbf{\_\_\_\_} & --- \\
\textbf{Coherence Score} & \textbf{\_\_\_\_ ÷ 6} & --- \\
\midrule
\end{tabular}
\end{center}

\subsubsection*{4.2 Coherence Score Interpretation}

\begin{center}
\begin{tabular}{@{}lll@{}}
\toprule
\textbf{Coherence Score} & \textbf{Assessment} & \textbf{Action Required} \\
\midrule
> 0.35 & Excellent integration & Maintain; document as best practice \\
0.25–0.35 & Acceptable coherence & Monitor; improve weak pairs \\
0.10–0.25 & Below acceptable & Redesign policies; strengthen spillovers \\
< 0.10 & Poor coordination & Major revision needed \\
\midrule
\end{tabular}
\end{center}

\subsubsection*{4.3 Identifying Weak Links (Low γ Pairs)}

\textbf{For each pair with γ < 0.2, answer:}

\begin{quote}
\textbf{Example: Health ↔ Pension, estimated γ = 0.05}

\textbf{Problem:}
\begin{itemize}[nosep]
\item Health policy says: ``Invest in behavior change now'' (present focus)
\item Pension policy says: ``Work longer for financial security'' (future focus)
\item Message conflict: Population hears: ``Confused priorities; what should I do?''
\end{itemize}

\textbf{Solution Options:}
\begin{enumerate}
\item \textbf{Reframe:} Connect health to work capacity: ``Healthy lifestyle enables longer, more productive working life''
\item \textbf{Sequence:} Launch health campaign first (build awareness), then pension reform (capitalize on behavior-change mentality)
\item \textbf{Segment:} Target different segments: Health to younger (build habits), Pension to 45+ (plan future)
\item \textbf{Accept:} If spillover γ ≈ 0 anyway, no interference; policies are independent
\end{enumerate}
\end{quote}

---

% =============================================================================
% SECTION 5: CONFLICT DETECTION DECISION TREE
% =============================================================================

\subsection*{Section 5: Policy Conflict Decision Tree}

\textbf{Use this tree to identify and prioritize conflicts:}

\begin{quote}
\textbf{START: Do you have policies in ≥2 domains? YES → Continue}

\textbf{STEP 1:} Compute γ for all domain pairs

→ Any γ < -0.2? \textbf{YES:} CONFLICT RISK! Go to Step 2
→ All γ ≥ -0.2? \textbf{NO:} Proceed to Step 3

\textbf{STEP 2 (Conflict Response):}
\begin{itemize}[nosep]
\item What policies create the negative spillover?
\item Are they both essential? Or can one be delayed/modified?
\item Option A: Redesign one policy to be compatible
\item Option B: Sequence policies (time-separate to reduce interference)
\item Option C: Communicate explicitly about the tradeoff
\end{itemize}

Go back to Step 1 after redesign.

\textbf{STEP 3 (Optimization):} Is Coherence Score > 0.25?

→ YES: Acceptable! Monitor quarterly.
→ NO: Identify weakest 2-3 pairs (lowest γ).

\textbf{STEP 4 (Strengthen Weak Pairs):}
For each pair γ < 0.2:
\begin{itemize}[nosep]
\item Can you add a bridging policy (e.g., job training to link Pension and Labor)?
\item Can you align messaging across domains?
\item Can you sequence policies better?
\end{itemize}

Recompute Coherence Score. Target ≥ 0.25.
\end{quote}

---

% =============================================================================
% SECTION 6: CASE STUDIES - POLICY CONFLICTS AND SOLUTIONS
% =============================================================================

\subsection*{Section 6: Case Studies of Policy Conflicts and Resolutions}

\subsubsection*{6.1 CASE STUDY 1: The Pension-Health Trap}

\textbf{Context:}
\begin{quote}
A government simultaneously launched:
\begin{itemize}[nosep]
\item \textbf{Health initiative:} ``Reduce obesity through behavior change and physical activity''
\item \textbf{Pension reform:} ``Work longer to secure retirement (age 67 instead of 65)''
\end{itemize}

Initial spillover estimate: γ(Health, Pension) = +0.05 (weak positive)
\end{quote}

\textbf{Problem Discovered:}
\begin{quote}
Population confused by contradictory urgencies:
\begin{itemize}[nosep]
\item Health campaign: ``Change your lifestyle right now---your future health depends on it!''
\item Pension reform: ``We're raising retirement age because we need you working 2 years longer''
\end{itemize}

Interpretation: ``The government wants me to invest in health improvements, but then immediately asks me to work longer despite aging. Which is it?''

Result: Both initiatives underperformed (50\% uptake vs. 70\% target)
\end{quote}

\textbf{Solution Applied:}
\begin{quote}
\textbf{Reframed} pension message: ``Working longer is feasible and rewarding if you invest in your health now. Healthy lifestyle means more energy, fewer sick days, and ability to continue doing what you love into your later years.''

Combined messaging: ``Health investment enables longer careers. Longer careers enable secure retirement.''

Result: Spillover improved γ = +0.25 (weak positive → moderate positive)
Uptake: Both initiatives reached 78-82\% (vs. earlier 50\%)
\end{quote}

\subsubsection*{6.2 CASE STUDY 2: The Environment-Labor Transition}

\textbf{Context:}
\begin{quote}
A city implemented:
\begin{itemize}[nosep]
\item \textbf{Environment:} ``Close coal power plants; transition to renewables''
\item \textbf{Labor:} ``Reskill coal workers for green energy sector''
\end{itemize}

Estimated spillover: γ(Environment, Labor) = +0.2 (weak positive)
\end{quote}

\textbf{Problem:}
\begin{quote}
Environment campaign emphasized job losses: ``Coal is dying; we're replacing it with renewables.'' This created fear among workers.

Labor initiative tried to reframe: ``New green jobs are coming!'' But workers heard this as too-good-to-be-true, viewed it as government spin.

Result: Labor unrest, strikes, delay of both initiatives
\end{quote}

\textbf{Solution:}
\begin{quote}
\textbf{Sequencing + messaging redesign:}
\begin{enumerate}
\item Month 0-6: Launch labor initiative FIRST (``Green energy jobs are growing; here's how to train'')
\item Month 6-12: Then announce environment policy: ``We're transitioning because green energy is already creating jobs in our region''
\item Messaging: ``Building on our success in green jobs, we're accelerating the transition''
\end{enumerate}

Result: Spillover improved γ = +0.4 (weak → strong positive)
Union support increased from 20\% to 65\% vs. original proposal
\end{quote}

---

% =============================================================================
% SECTION 7: MULTI-DOMAIN GOVERNANCE CHECKLIST
% =============================================================================

\subsection*{Section 7: Multi-Domain Policy Governance Checklist}

\textbf{Use this before implementing any cross-domain policy package:}

\begin{center}
\small
\begin{tabular}{@{}p{1cm}p{8cm}p{1.5cm}@{}}
\toprule
\textbf{\#} & \textbf{Checklist Item} & \textbf{Status} \\
\midrule
1 & All domain policies identified and documented & ☐ \\
2 & Spillover matrix completed for all 6 pairs & ☐ \\
3 & Coherence Score computed (target >0.25) & ☐ \\
4 & Any γ < -0.2 identified and conflict mitigation planned & ☐ \\
5 & Weak pairs (γ < 0.2) have strengthening strategy & ☐ \\
6 & Messaging across domains checked for consistency & ☐ \\
7 & Segment targeting verified (no contradictions) & ☐ \\
8 & Phase progression sequencing planned (if needed) & ☐ \\
9 & Implementation timeline coordinate across domains & ☐ \\
10 & Quarterly review schedule set up to monitor spillovers & ☐ \\
\midrule
\end{tabular}
\end{center}

---

% =============================================================================
% SECTION 8: QUARTERLY CROSS-DOMAIN REVIEW PROTOCOL
% =============================================================================

\subsection*{Section 8: Quarterly Cross-Domain Spillover Review}

\textbf{Every quarter, review this protocol to detect emerging conflicts:}

\begin{center}
\begin{tabular}{@{}llccc@{}}
\toprule
\textbf{Metric} & \textbf{Definition} & \textbf{Q\_\_} & \textbf{Q+1} & \textbf{Trend} \\
\midrule
\textbf{Coherence} & Sum of γ ÷ 6 & \_\_\_\_ & \_\_\_\_ & ↑ / → / ↓ \\
\textbf{Weak Pairs} & Count of γ < 0.2 & \_\_\_\_ & \_\_\_\_ & ↑ / → / ↓ \\
\textbf{Conflicts} & Count of γ < -0.2 & \_\_\_\_ & \_\_\_\_ & ↑ / → / ↓ \\
\textbf{Message Clarity} & Population confusion score & \_\_\_\_ & \_\_\_\_ & ↑ / → / ↓ \\
\multicolumn{5}{l}{\small↑ = Improving, → = Stable, ↓ = Deteriorating} \\
\midrule
\end{tabular}
\end{center}

\textbf{If Coherence Score drops below 0.25:}
\begin{itemize}[nosep]
\item Identify which pair deteriorated
\item Conduct root cause analysis (message changed? implementation problem?)
\item Redesign affected policy or messaging
\item Re-estimate γ and recompute Coherence Score
\end{itemize}

---

% =============================================================================
% REFERENCES
% =============================================================================

\subsection*{References and Related Reading}

\textbf{For Policy Designers:}
\begin{itemize}[nosep]
\item Chapter 22: Policy Applications (shows domains with examples)
\item Appendix LLL: METHOD-POLICY-DOMAINS (domain-specific templates)
\end{itemize}

\textbf{For Theory:}
\begin{itemize}[nosep]
\item Appendix HHH (METHOD-TOOLKIT): Coherence criteria C1-C4
\item Appendix KKK (METHOD-DYNAMIC): Formal spillover dynamics
\end{itemize}

\textbf{For Measurement:}
\begin{itemize}[nosep]
\item Section 8: Quarterly review protocol
\item Appendix R: Evaluation Protocol (for each domain separately)
\end{itemize}


%%%%%%%%%%%%%%%%%%%%%%%%%%%%%%%%%%%%%%%%%%%%%%%%%%%%%%%%%%%%%%%%%%%%%%%%%%%%%%%

%%%%%%%%%%%%%%%%%%%%%%%%%%%%%%%%%%%%%%%%%%%%%%%%%%%%%%%%%%%%%%%%%%%%%%%%%%%%%%%
%%%%%%%%%%%%%%%%%%%%%%%%%%%%%%%%%%%%%%%%%%%%%%%%%%%%%%%%%%%%%%%%%%%%%%%%%%%%%%%
\section{The Fundamental Question}
\label{sec:mmm-fundamental}
%%%%%%%%%%%%%%%%%%%%%%%%%%%%%%%%%%%%%%%%%%%%%%%%%%%%%%%%%%%%%%%%%%%%%%%%%%%%%%%

\begin{tcolorbox}[
    colback=red!5!white,
    colframe=red!75!black,
    title={\textbf{Fundamental Question}}
]
\textbf{How does unknown integrate with and contribute to the
Evidence-Based Framework for understanding behavioral outcomes?}

This appendix addresses the core challenge of formalizing behavioral principles
within a rigorous theoretical framework while maintaining practical applicability.
\end{tcolorbox}


\section{Critical Foundations and Objections}
\label{sec:mmm-foundations}
%%%%%%%%%%%%%%%%%%%%%%%%%%%%%%%%%%%%%%%%%%%%%%%%%%%%%%%%%%%%%%%%%%%%%%%%%%%%%%%

\subsection{Objection 1: Scope Limitations}

\textbf{Objection:} The framework presented may have limited applicability
outside the specific domain contexts discussed.

\textbf{Response:} We acknowledge domain-specificity. The framework provides
general principles that should be calibrated to specific contexts. Cross-domain
validation remains an open research question.

\subsection{Objection 2: Measurement Challenges}

\textbf{Objection:} Key constructs may be difficult to measure reliably in
practice.

\textbf{Response:} We recommend using validated scales and instruments where
available, and developing domain-specific proxies where necessary. Sensitivity
analysis should assess robustness to measurement uncertainty.



%%%%%%%%%%%%%%%%%%%%%%%%%%%%%%%%%%%%%%%%%%%%%%%%%%%%%%%%%%%%%%%%%%%%%%%%%%%%%%%
%%%%%%%%%%%%%%%%%%%%%%%%%%%%%%%%%%%%%%%%%%%%%%%%%%%%%%%%%%%%%%%%%%%%%%%%%%%%%%%

%%%%%%%%%%%%%%%%%%%%%%%%%%%%%%%%%%%%%%%%%%%%%%%%%%%%%%%%%%%%%%%%%%%%%%%%%%%%%%%
\section{Key Results}
\label{sec:mmm-results}
%%%%%%%%%%%%%%%%%%%%%%%%%%%%%%%%%%%%%%%%%%%%%%%%%%%%%%%%%%%%%%%%%%%%%%%%%%%%%%%

The analysis yields the following key results:

\begin{enumerate}
    \item \textbf{Result 1:} [Primary finding with quantitative support]
    \item \textbf{Result 2:} [Secondary finding with implications]
    \item \textbf{Result 3:} [Practical application or boundary condition]
\end{enumerate}


\section{Summary}
\label{sec:mmm-summary}
%%%%%%%%%%%%%%%%%%%%%%%%%%%%%%%%%%%%%%%%%%%%%%%%%%%%%%%%%%%%%%%%%%%%%%%%%%%%%%%

This appendix has presented the key concepts and theoretical foundations.
The main contributions include:

\begin{enumerate}
    \item Formal definitions and theoretical framework
    \item Integration with the broader EBF structure
    \item Practical guidelines for application
\end{enumerate}

The content supports decision-making and behavioral analysis within the
Evidence-Based Framework, providing a foundation for further research
and practical implementation.


\section{Glossary of Symbols}
\label{sec:mmm-glossary}
%%%%%%%%%%%%%%%%%%%%%%%%%%%%%%%%%%%%%%%%%%%%%%%%%%%%%%%%%%%%%%%%%%%%%%%%%%%%%%%

For comprehensive definitions, see Appendix G (Master Glossary).

\begin{table}[htbp]
\centering
\caption{Key Symbols in Appendix MMM}
\begin{tabular}{lll}
\toprule
\textbf{Symbol} & \textbf{Meaning} & \textbf{Reference} \\
\midrule
$\gamma$ & Complementarity parameter & Appendix B \\
$\Psi$ & Context vector & Appendix V \\
$U$ & Utility function & Chapter 3 \\
\bottomrule
\end{tabular}
\end{table}


\section{References}
\label{sec:references}
%%%%%%%%%%%%%%%%%%%%%%%%%%%%%%%%%%%%%%%%%%%%%%%%%%%%%%%%%%%%%%%%%%%%%%%%%%%%%%%

See master bibliography (\texttt{bcm\_master.bib}) for complete citations.

\nocite{bcm_master}

% =============================================================================
% END OF APPENDIX MMM
% =============================================================================

\vspace{2cm}
\begin{center}
{\large \textbf{End of Appendix MMM: METHOD-POLICY-COHERENCE}} \\
\vspace{0.5cm}
{\small \textit{For domain-specific design templates, see Appendix LLL}}
\end{center}
